\section*{摘要}

随着2022年FIA引入"地面效应"空气动力学规则，
F1赛车的海豚跳（Porpoising）现象成为工程设计中的关键挑战。
传统CFD仿真在规则限制下难以高效覆盖大规模参数空间，
亟需低成本、高效率的替代方法。

本文基于Kaggle公开的F1气动稳定性数据集（约15万样本），
该数据集由XGBoost预测引擎在NVIDIA P100上生成（生成精度99.66\%），
构建了一套数据驱动的代理优化框架。
首先通过6种机器学习模型的系统对比，
基于实验证据将技术路线从开题的"BP神经网络"调整为
XGBoost（可解释性）与DeepEnsemble（PSO寻优）的双模型策略，
分别实现R$^2 = 0.9999$和R$^2 = 0.9995$的预测精度，
并通过5折交叉验证排除了数据泄漏的可能。
在此基础上，
将不确定性感知（Uncertainty-Aware）思想引入PSO框架，
以$\mu(x)-\lambda\sigma(x)$为适应度函数，
有效抑制代理模型在分布外（OOD）区域的幻觉行为。
进一步通过不确定性引导的模型精炼机制（Uncertainty-Guided Refinement），
将MSE从0.3272降至0.1444（降低55.9\%），
收敛速度提升2.1倍。

在多场景优化模块中，
针对Monza高速、Monaco高下压力、均衡设置和湿地四种典型F1工况，
采用多目标加权适应度函数进行Pareto前沿搜索，
并通过SHAP、排列重要性、约束灵敏度分析和Landscape Diagnosis
对最优解进行了系统性的可解释性验证。
Porpoising风险分析模块通过一阶梯度热力图和二阶Hessian曲率分析，
刻画了不同工况下的海豚跳风险空间分布。

实验结果表明：
代理模型R$^2$达到0.9999，远超理想目标（$\geq 0.92$）；
PSO平均收敛代数为28.9代（最优目标$\leq 50$）；
单次寻优耗时约72ms，计算效率优于GA（970ms）和DE（769ms）；
消融实验验证了不确定性感知惩罚的OOD保护功能，
以及多目标优化相较单目标优化的区分度提升（从0.1\%到23\%以上）。
景观诊断揭示了目标函数退化现象——
86\%样本中stability趋近满分导致优化景观平坦，
但这一退化源自数据分布，而非PSO或代理模型的缺陷。

研究工作形成了一套完整的"数据预处理—代理建模—不确定性感知优化—可解释性验证—景观诊断"
技术路线，
为代理模型驱动的工程优化提供了可量化的方法论支撑。
本研究的核心发现是：在数据驱动的优化中，
优化器选择的影响远小于数据分布结构本身——
当86\%的样本趋近饱和时，PSO、GA、DE的结果趋于一致，
优化景观而非优化算法决定了寻优上限。
景观诊断方法本身可迁移至任何使用代理模型进行优化的工程场景。

\vspace{0.5cm}

\noindent
\textbf{关键词：}
代理模型；
粒子群优化；
不确定性感知优化；
计算智能；
F1赛车；
气动稳定性；
海豚跳；
多目标优化；
景观诊断；
可解释性分析
